% Technical summary of implementation results — for inclusion in the final
% report. Same register and conventions as the project proposal.
\documentclass[12pt,a4paper]{article}
\usepackage[a4paper,margin=1in]{geometry}
\usepackage{amsmath}
\usepackage{booktabs}
\usepackage{setspace}
\usepackage[colorlinks=true, linkcolor=red, citecolor=green!60!black]{hyperref}
\setstretch{1.25}
\setlength{\parindent}{0pt}
\setlength{\parskip}{7pt}

\title{\vspace{-1.5cm}\Large\bfseries LiDAR Point Cloud Denoising for Autonomous
Driving Applications\\[4pt]\large Technical Summary of Implementation Results}
\author{CPG No.~5 --- Department of Computer Science and Engineering,\\
Thapar Institute of Engineering and Technology, Patiala}
\date{July 2026}

\begin{document}
\maketitle
\vspace{-0.8cm}

\section*{1. System Implemented}

The pipeline proposed in Section~5 of the project proposal has been
implemented in full. Raw KITTI Velodyne scans are preprocessed (NaN/Inf
removal, 60~m range gating, 0.15~m voxel-grid downsampling), corrupted with
the three noise families identified in the Need Analysis (Gaussian jitter,
$\sigma=0.05$~m; uniform outliers, 8\%; fog-style near-range backscatter,
4\%, exponentially distributed range), and denoised by \emph{both} approaches
named in the methodology: (a) a statistical baseline chaining Statistical
Outlier Removal ($k=20$, $\alpha=2.0$) and Radius Outlier Removal (4
neighbours within 0.6~m), and (b) two learning-based models trained on
paired noisy/clean patches --- a PointNet autoencoder (0.83M parameters,
Chamfer loss) and a lightweight DGCNN (0.22M parameters) that regresses
per-point displacements and classifies injected noise points for removal.

\textbf{Data.} 262 scans from two public KITTI raw drives
(2011\_09\_26, drives 0001 and 0005), split 213/26/28 into
train/validation/test by frame. The official KITTI 3D-object Velodyne
archive (29~GB, login-gated) was deliberately not used: the raw split
provides identical sensor data at 490~MB, sufficient for denoising
experiments. Training used $\sim$39{,}000 patches of 256 points; the largest
run feasible on team hardware (Apple M1, MPS backend) was 30 epochs for the
autoencoder and 10 for the DGCNN.

\section*{2. Reconstruction Quality (Proposal Table-1 axes)}

Averages over 10 held-out test frames; CD = Chamfer Distance,
HD$_{95}$ = 95th-percentile Hausdorff (robust to single stragglers).

\begin{table}[h]
\centering
\begin{tabular}{lccccc}
\toprule
\textbf{Method} & \textbf{CD (m)} & \textbf{HD$_{95}$ (m)} &
\textbf{CD impr.} & \textbf{Retention} & \textbf{Latency}\\
\midrule
Noisy input     & 0.403 & ---   & ---    & ---    & ---\\
SOR + ROR       & \textbf{0.095} & \textbf{0.136} & \textbf{75.5\%} & 88.6\% & \textbf{104 ms}\\
PointNet AE     & 0.374 & 1.835 & 5.4\%  & 165.6\% & 282 ms\\
DGCNN           & 0.162 & 0.156 & 59.1\% & 92.3\% & 3{,}011 ms\\
\bottomrule
\end{tabular}
\end{table}

The statistical baseline remains the strongest pure-reconstruction method,
consistent with the literature survey's expectation that classical filters
are highly effective against \emph{additive outlier} noise. The DGCNN
approaches it (and matches it on HD$_{95}$) while learning the noise model
rather than assuming it. The autoencoder under-performs: reconstructing
patches through a global bottleneck over-smooths fine scan-line structure
and inflates the cloud (retention $>$100\%), a known limitation of
set-decoder architectures. An ablation over noise intensity
($\sigma = 0.02$--$0.12$~m) shows SOR+ROR and DGCNN degrading gracefully
and in parallel, while the AE's error is dominated by reconstruction blur
and is insensitive to $\sigma$. Sweeping SOR parameters confirms the
proposal's tuning discussion: tighter thresholds ($\alpha=1.0$) improve CD
by $\sim$6\% over the default on this noise profile.

\section*{3. Downstream Detection Impact}

KITTI raw frames carry no ground-truth boxes, so a \emph{detection
consistency} protocol was used: a KITTI-pretrained PointPillars detector
(MMDetection3D, CPU inference) is run on each clean cloud and its boxes
taken as reference; recall/precision of the detections on noisy and denoised
clouds are measured against that reference (BEV centre matching, 1~m
tolerance, score $\geq 0.4$, 10 frames).

\begin{table}[h]
\centering
\begin{tabular}{lcc}
\toprule
\textbf{Condition} & \textbf{Consistency recall} & \textbf{Consistency precision}\\
\midrule
Noisy input & 0.760 & 0.853\\
SOR + ROR   & 0.733 & \textbf{0.975}\\
PointNet AE & 0.067 & 0.750\\
DGCNN       & \textbf{0.788} & 0.870\\
\bottomrule
\end{tabular}
\end{table}

The synthetic noise costs the detector a quarter of its detections and
introduces spurious boxes. The two viable denoisers repair different
failure modes: SOR+ROR almost eliminates noise-induced \emph{false} boxes
(precision 0.85$\to$0.98) at a small recall cost, while the DGCNN recovers
the most \emph{true} detections (recall 0.79, the best of all conditions).
Notably, the ranking by Chamfer Distance does not predict this: the AE, only
7\% worse than noisy on CD, is catastrophic downstream (recall 0.07). This
directly substantiates the research gap identified in Section~3.5 of the
proposal --- reconstruction metrics alone are an incomplete proxy for
perception quality.

\section*{4. Real-Time Feasibility (10 Hz target)}

Per-frame latency was profiled end-to-end on the Apple M1 (batch inference,
patch covering, k-d tree construction included). Dynamic int8 quantisation
was applied to the learning models as the Step-5 optimisation.

\begin{center}
\begin{tabular}{lrr}
\toprule
\textbf{Configuration} & \textbf{ms/frame} & \textbf{Hz}\\
\midrule
SOR+ROR (CPU)                & 174    & 5.7\\
PointNet AE (MPS GPU)        & 138    & 7.3\\
PointNet AE (CPU fp32)       & 2{,}856 & 0.35\\
PointNet AE (CPU int8)       & 2{,}898 & 0.35\\
DGCNN (MPS GPU)              & 3{,}317 & 0.30\\
DGCNN (CPU fp32)             & 20{,}437 & 0.05\\
\bottomrule
\end{tabular}
\end{center}

\textbf{Honest assessment: the 10~Hz target is not yet met.} The closest
configurations are the GPU-batched autoencoder (7.3~Hz) and the statistical
baseline (5.7~Hz, dominated by ROR's radius queries); the DGCNN's dynamic
k-NN graph construction is the principal bottleneck (0.3~Hz). Quantisation
reduced the AE checkpoint 3.4$\times$ (3.34$\to$0.98~MB) --- valuable for
the proposal's embedded Raspberry Pi scenario --- but did not improve
latency, as PyTorch's qnnpack dynamic backend accelerates only the Linear
layers, which are not the bottleneck. The measured path to 10~Hz, left as
future work, is (i)~SOR-only filtering (interior 90~ms of the 174~ms
pipeline), (ii)~precomputed/voxel-hashed neighbourhoods for the DGCNN, and
(iii)~CUDA-class hardware, on which the batched models have substantial
headroom.

\section*{5. Deliverables and Scope Notes}

The repository contains the modular pipeline
(\texttt{src/lidar\_denoise/}), a YAML-configured driver
(\texttt{scripts/run\_pipeline.py}), the detection evaluation
(\texttt{scripts/detection\_eval.py}), an interactive before/after demo
(\texttt{scripts/demo.py}), 14 unit tests, all metrics/plots under
\texttt{results/}, and the corrected proposal source. The YDLIDAR~X2 /
Raspberry~Pi live-capture component of the proposal is hardware-dependent
and deferred; the pipeline is sensor-agnostic and accepts any $(N,3)$ scan.

\end{document}
